\newpage
\phantomsection
\label{prf:ec-addition}

\begin{remark*}[Return]
\hyperref[thm:ec-addition]{Return to Theorem}
\end{remark*}

\begin{theorem*}[Addition of Equivalence Classes]
For every $[(a,b)],[(c,d)]\in\mathbb{Q}$,
\[
[(a,b)] + [(c,d)] = [(ad+bc,\,bd)].
\]
\end{theorem*}

\begin{proof}
\textbf{Professional Standard Proof.}~\\
Let $[(a,b)],[(c,d)]\in\mathbb{Q}$. Then $a,b,c,d\in\mathbb{Z}$ with
$b\neq 0$ and $d\neq 0$.

By the definition of addition on rational equivalence classes,
\[
[(a,b)] + [(c,d)]
:=
[(ad+bc,bd)].
\]
Therefore
\[
[(a,b)] + [(c,d)] = [(ad+bc,bd)].
\]
The well-definedness lemma guarantees that this value is independent of
the representatives chosen for the two equivalence classes.
\end{proof}

\begin{proof}
\textbf{Detailed Learning Proof.}~\\

\textbf{Step 1.} Interpret the hypotheses as rational equivalence classes.

Let
\[
[(a,b)],[(c,d)]\in\mathbb{Q}.
\]
Then $(a,b)$ and $(c,d)$ are admissible fraction-pairs. Hence
\[
a,b,c,d\in\mathbb{Z}
\qquad\text{with}\qquad
b\neq 0
\qquad\text{and}\qquad
d\neq 0.
\]

\textbf{Step 2.} Apply the definition of addition on rational classes.

Addition on rational equivalence classes is defined by the rule
\[
[(a,b)] + [(c,d)]
:=
[(ad+bc,bd)].
\]
Therefore, directly from the definition,
\[
[(a,b)] + [(c,d)] = [(ad+bc,bd)].
\]

\textbf{Step 3.} Record why the formula is legitimate.

The expression on the right is independent of the chosen representatives.
Indeed, this is exactly the content of the well-definedness lemma for
rational addition. Thus the displayed formula defines an operation on
equivalence classes, not merely on representatives.
\end{proof}

\begin{remark*}[Proof structure]
The proof is definitional. Rational addition has already been defined on
equivalence classes by the displayed representative formula. The
well-definedness lemma supplies the prior justification that the formula
does not depend on the choice of representatives.
\end{remark*}

\begin{remark*}[Dependencies]~\\
\begin{itemize}
  \item \hyperref[def:rational-operations]{Definition~\ref*{def:rational-operations}}
        \textnormal{(addition defined by the formula; this theorem restates it as an explicit identity)}
  \item \hyperref[lem:rational-addition-well-defined]{Lemma~\ref*{lem:rational-addition-well-defined}}
        \textnormal{(well-definedness: result is independent of representative choice)}
\end{itemize}
\end{remark*}